\documentclass[mathserif]{beamer}
\usepackage{subfigure}
\usepackage{listings}
\usepackage{courier}
\definecolor{lightlightgray}{gray}{0.95}
\definecolor{lightlightblue}{rgb}{0.4,0.4,0.95}
\definecolor{lightlightgreen}{rgb}{0.8,1,0.8}
\lstset{language=C++,
           frame=single,
           basicstyle=\ttfamily\footnotesize,
           keywordstyle=\color{black}\textbf,
           backgroundcolor=\color{lightlightgray},
           commentstyle=\color{blue},
           frame=single
           }
\usetheme[secheader]{pecostalk}
\usepackage{bibentry}
\nobibliography*
\graphicspath{{figs/}}                                                                                                                              

% defines newblock as null, giving compile issues otherwise
\let\newblock\relax 

\newcommand{\eqdef}{\stackrel{\text{\tiny def}}{=}}

\newcommand{\NVRvect}[1]{\ensuremath\boldsymbol{#1}}
\newcommand{\vect}[1]{\ensuremath\boldsymbol{#1}}
\newcommand{\NVRtensor}[1]{\NVRvect{#1}}
%\newcommand{\NVRtensor}[1]{\underline{\NVRvect{#1}}}
\newcommand{\NVRnorm}[1]{\left|\left|#1\right|\right|}
\newcommand{\norm}[1]{\left|\left|#1\right|\right|}
\newcommand{\NVRgrad}{\nabla}
\newcommand{\NVRdiv}{\NVRgrad \cdot}
\newcommand{\NVRpd}[2]{\frac{\partial#1}{\partial#2}}
\newcommand{\NVRpdd}[2]{\frac{\partial^2#1}{\partial#2^2}}
\newcommand{\NVReqdef}{\stackrel{\text{\tiny def}}{=}}

\newcommand{\NVRHgrad}{H(\text{grad})}
\newcommand{\NVRHdiv}{H(\text{div})}
\newcommand{\NVRsumm}[2]{\ensuremath\displaystyle\sum\limits_{#1}^{#2}}
 
\newcommand{\NVRcurl}{\nabla \times}
\newcommand{\NVRHcurl}{\ensuremath H(\text{curl})}

\newcommand{\code}[1]{\texttt{#1}}
\newcommand{\deal}{\code{deal.II}\,}

\newcommand{\pecosbold}[1]{{\color{pecos2}{#1}}}
\newcommand{\pecosreallybold}[1]{{\color{pecos6}{#1}}}

\newcommand{\FootSize}{\scriptsize}

\DeclareMathOperator*{\argmin}{arg\,min}

\date{September 12, 2012}
\author[Nate Roberts]{Nate Roberts \\
Collaborators: Leszek Demkowicz, Jesse Chan, Truman Ellis
}
\institute{Institute for Computational and Engineering Sciences\\
The University of Texas at Austin}
\title[Camellia: A DPG Toolbox]{Camellia:}
\subtitle{A Toolbox for a Class of Discontinuous Petrov-Galerkin Methods\\ Using Trilinos}

\begin{document}
\begin{frame}
%\begin{center}
%\includegraphics[width=.8\linewidth]{grand_logo}\\
%\end{center}
\titlepage
%\begin{flushright}
%\includegraphics[scale=0.1]{asc_logo}\\
%\end{flushright}
\end{frame}

%===============================================================================
% OUTLINE
%===============================================================================

% UNCOMMENT THIS TO ADD AN OUTLINE AT THE TOP OF EACH SECTION...
%\AtBeginSection[]
%{
%  \begin{frame}<beamer>
%    \frametitle{Outline}
%    \tableofcontents[currentsection]
%  \end{frame}
%}

\begin{frame}
\frametitle{Outline}
\tableofcontents
\end{frame}

\section{DPG in Brief} % 5-6 slides
%===============================================================================
% From Strong-Form PDE to DPG Form
%===============================================================================
\begin{frame}                                                                                                                                                                          
\frametitle{From Strong-Form PDE to DPG Form}
\begin{center}
\includegraphics[width=\linewidth]{../figures/DPGFormCartoon}\\
\end{center}
\end{frame}              
%===============================================================================
% Solving with DPG
%===============================================================================
\begin{frame}                                                                                                                                                                          
\frametitle{Solving with DPG}
\begin{center} 
\includegraphics[width=0.9\linewidth]{../figures/DPGSolveCartoonNew}\\
\end{center}
\end{frame}

\begin{frame}
\frametitle{Graph Test Norm}

For a strong operator $A$ with formal adjoint $A^{*}$, the \pecosbold{adjoint graph space} is
\[
H_{A^{*}} = \left\{ v \in L^{2}(\Omega) : A^{*}v \in L^{2}(\Omega) \right\}
\]

and the \pecosbold{(adjoint) graph norm} on the test space $V$ is given by
\[
\NVRnorm{v}_{\rm graph} = \NVRnorm{v}_{H_{A^{*}}} = \left( \NVRnorm{v}^{2} + \NVRnorm{A^{*}v}^{2}\right)^{1/2}.
\]

\vspace{5mm}
\begin{center}
E.g. if $A^{*} = \NVRgrad$, then $H_{A^{*}} = H^{1}$, and $\NVRnorm{v}_{H_{A^{*}}} = \NVRnorm{v}_{H^{1}}$.
\end{center}
\end{frame}

\begin{frame}
\frametitle{Key Result: Well-posedness $\implies$ Optimal Convergence}
Under modest technical assumptions (true for Stokes), we have\footnote{\FootSize \bibentry{DPGStokes}}\\
\[
\NVRnorm{Au} \geq \gamma \NVRnorm{u} \implies \sup_{v \in H_{A^{*}}} \frac{b((u,\widehat{u}),v)}{\NVRnorm{v}_{H_{A^{*}}}} \geq \gamma_{\rm DPG} \left(\NVRnorm{u}^{2} + \NVRnorm{\widehat{u}}^{2}_{\widehat{H}_{A}(\Gamma_{h})}\right)^{1/2}
\]
where $\gamma_{\rm DPG} = O(\gamma)$ is a mesh-independent constant, and $\NVRnorm{\cdot}_{\widehat{H}_{A}(\Gamma_{h})}$ is the minimum energy extension norm.\\
\pause
\vspace{5mm}
By Babu\v{s}ka's Theorem,
\[
\implies \left(\NVRnorm{u - u_{h}}^{2} + \NVRnorm{\widehat{u} - \widehat{u}_{h}}^{2}_{\widehat{H}_{A}(\Gamma_{h})}\right)^{1/2} \leq \frac{M}{\gamma_{\rm DPG}} \left( {\rm B.A.E.} \right).
\]
\pause
Suffices to show that $\NVRnorm{Au} \geq \gamma \NVRnorm{u}$ to prove \pecosbold{optimal} convergence rate!
\end{frame}

\section{Stokes} % 10-12 slides
% lit. review (including various classical elements, discussion of locking and non-convergence)
% Stokes analysis, summary results
% Numerical experiments: manufactured solution, lid-driven cavity flow problem
% Probably worth mentioning the naive norm results: drives home why the analysis is valuable, and the fact that the optimal test function lives in a very specific space.
%===============================================================================
% Stokes: Classical Problem
%===============================================================================
\begin{frame}
\frametitle{Classical Stokes Problem}
The classical strong form of the Stokes problem in $\Omega \subset
\mathbb{R}^{2}$ is given by
\begin{align*}
- \mu \Delta \vect{u} + \NVRgrad p &= \vect{f} & \text{ in } \Omega, \\
\NVRdiv \vect{u} &= 0 & \text{ in } \Omega, \\
\vect{u} &= \vect{u}_D & \text{ on } \partial\Omega,
\end{align*}
where $\mu$ is viscosity, $p$ pressure, $\vect{u}$ velocity, and
$\vect{f}$ a vector forcing function.  Since by appropriate
non-dimensionalization we can eliminate the constant $\mu$, we take
$\mu=1$ throughout. 
\end{frame}

\begin{frame}
\frametitle{DPG Applied to Stokes}

To apply DPG, we need a first-order system.  We introduce $\NVRtensor{\sigma}= \NVRgrad
\vect{u}$:
\begin{align*}
- \NVRdiv \NVRtensor{\sigma} + \NVRgrad p &= \vect{f} & \text{ in } \, \Omega, \\
 \NVRdiv \vect{u} &= 0 & \text{ in } \, \Omega, \\
 \NVRtensor{\sigma} - \NVRgrad \vect{u} &= 0 & \text{ in } \, \Omega.
\end{align*}
Testing with $(\vect{v},q,\vect{\tau})$, and integrating by parts, we have
\begin{align*}
\left(\NVRtensor{\sigma} - p \NVRtensor{I}, \NVRgrad \vect{v} \right)_{\Omega_{h}} - \left\langle \widehat{\vect{t}}_{n}, \vect{v} \right\rangle_{\Gamma_{h}} &= (\vect{f}, \vect{v} )_{\Omega_{h}}\\
\left(\vect{u}, \NVRgrad q \right)_{\Omega_{h}} -  \left\langle \widehat{\vect{u}} \cdot \vect{n}, q  \right\rangle_{\Gamma_{h}} &= 0\\
\left( \NVRtensor{\sigma}, \NVRtensor{\tau} \right)_{\Omega_{h}} + \left( \vect{u}, \NVRdiv \NVRtensor{\tau} \right)_{\Omega_{h}} - \left\langle \widehat{\vect{u}}, \NVRtensor{\tau} \vect{n}  \right\rangle_{\Gamma_{h}} &= \NVRtensor{0},
\end{align*}
where traction $\vect{t}_{n} \eqdef ( \NVRtensor{\sigma} - p \NVRtensor{I} ) \vect{n}$, and the hatted variables $\widehat{\vect{t}}_{n}$ and $\widehat{\vect{u}}$ are new unknowns representing the traces of the corresponding variables at the boundary.

\end{frame}

\begin{frame}
\frametitle{DPG Applied to Stokes}
DPG Formulation:
\begin{align*}
b(u,v) =& \left(\NVRtensor{\sigma} - p \NVRtensor{I}, \NVRgrad \vect{v} \right)_{\Omega_{h}} - \left\langle \widehat{\vect{t}}_{n}, \vect{v} \right\rangle_{\Gamma_{h}}\\
&+\left(\vect{u}, \NVRgrad q \right)_{\Omega_{h}} -  \left\langle \widehat{\vect{u}} \cdot \vect{n}, q  \right\rangle_{\Gamma_{h}}\\
&+ \left( \NVRtensor{\sigma}, \NVRtensor{\tau} \right)_{\Omega_{h}} + \left( \vect{u}, \NVRdiv \NVRtensor{\tau} \right)_{\Omega_{h}} - \left\langle \widehat{\vect{u}}, \NVRtensor{\tau} \vect{n}  \right\rangle_{\Gamma_{h}} = (\vect{f}, \vect{v} )_{\Omega_{h}} = l(v).
\end{align*}
The natural spaces for the trial variables are then:
\begin{itemize}
\item fields: $p \in L^{2}(\Omega), \vect{u} \in \vect{L}^{2}(\Omega), \NVRtensor{\sigma} \in \vect{L}^{2}(\Omega)$,
\item fluxes: $\widehat{\vect{t}}_{n} \in \NVRtensor{H}^{-1/2}(\Gamma_{h})$,
\item traces: $\widehat{\vect{u}} \in \vect{H}^{1/2}(\Gamma_{h})$.
\end{itemize}

The natural norms for fluxes and traces are \emph{minimum energy extension norms}.

\end{frame}

%\begin{frame}
%\frametitle{Polynomial Space Choices}
%How to select polynomial spaces for $u$ and $\widehat{u}$?
%\begin{itemize}
%\item Graph norm $\implies$ error bounded by best approximation error.
%\vspace{3mm} \pause
%\item Want all variables to have same B.A.E. convergence rates.
%\vspace{3mm} \pause
%\item Fix polynomial order $k$ for $H^{1}(\Omega_{h})$ variables.  ($k$ rate.)
%\vspace{3mm} \pause
%\item $\implies$ need order $k-1$ for $L^{2}(\Omega)$.
%\vspace{3mm} \pause
%\item $\implies$ need order $k$ for $H^{1/2}(\Gamma_{h})$.
%\vspace{3mm} \pause
%\item $\implies$ need order $k-1$ for $H^{-1/2}(\Gamma_{h})$.
%\vspace{3mm} \pause
%\item In Camellia, specify the $H^{1}$ polynomial order $k$, and the other orders will be chosen automatically.
%\end{itemize}
%\end{frame}

\begin{frame}
\frametitle{Graph Test Norm}
The adjoint graph norm for our Stokes formulation is:\footnote{\FootSize \bibentry{DPGStokes}}
\begin{align*}
\norm{(\NVRtensor{\tau},\vect{v}, q)}_{\rm graph}^{2} =& \norm{\NVRdiv \NVRtensor{\tau} - \NVRgrad q}^{2} + \norm{\NVRdiv \vect{v}}^{2} + \norm{ \NVRtensor{\tau} + \NVRgrad{\vect{v}}}^{2}\\ 
&+ \norm{\NVRtensor{\tau}}^{2} + \norm{\vect{v}}^{2} + \norm{q}^{2}.
\end{align*}

\end{frame}

\section{Camellia} % 10-12 slides
\begin{frame}
\frametitle{Camellia\footnote{\FootSize \bibentry{RobertsetAl11}}}
\pecosbold{Design Goal}: make DPG research and experimentation as simple as possible, without sacrificing too much by way of performance.
\vspace{3mm}
\begin{center}
\includegraphics[scale=.65]{../figures/CamelliaVennDiagram.pdf}
\end{center}
\addtocounter{footnote}{1}
\footnotetext{\FootSize \bibentry{TrilinosShortAuthors}}
\end{frame}

%===============================================================================
% WHAT INTREPID (TRILINOS) PROVIDES
%===============================================================================
\begin{frame}
\frametitle{Reliance on Trilinos}
\begin{block}{}
\begin{center}
\begin{tabular}{l c}
\pecosbold{Feature}			&\pecosbold{Trilinos Package} \\
OO interface to MUMPS 		&Amesos \\
KLU solver 				&Amesos \\
conforming basis functions	&Intrepid \\
pullbacks/Piola transforms 	&Intrepid \\
smart multidimensional arrays 	&Intrepid \\
distributed compressed row storage matrices	&Epetra \\
cell topologies 				&Shards \\
reference-counted pointers 				&Teuchos \\
space-filling curves for spatially local mesh partitioning 		&Zoltan \\
\end{tabular}
\end{center}
\end{block}
\end{frame}

\begin{frame}
\frametitle{Camellia: Users}
\begin{itemize}
\item Jesse Chan: compressible Navier-Stokes
\item Truman Ellis: compressible Navier-Stokes with turbulence
\item Nate Roberts: incompressible Navier-Stokes
\end{itemize}
\begin{figure}
\centering
\subfigure[Mesh]{\includegraphics[height = 3cm]{../figures/Jesse/Mesh10.png}}
\subfigure[$u_1$]{\includegraphics[height = 3cm]{../figures/Jesse/u10.png}}
\caption{Compressible Navier-Stokes for Carter flat plate problem, ${\rm Re}=1000, {\rm Ma}=3.$}
\end{figure}
\end{frame}

\section{Solving Stokes with Camellia}
\begin{frame}
\frametitle{Implementation with Camellia}
%PROBABLY BETTER TO MAKE A CARTOON HERE
\begin{enumerate}
\item Identify variables involved: fields, traces, fluxes, and test functions.
\item Define the bilinear form.
\item Define right-hand side.
\item Define inner product.
\item Define boundary conditions.
\item Create initial mesh.
\item Create solution object (and solve).
\item Adaptively refine the mesh.
\end{enumerate}
\end{frame}

\subsection{Identify variables}
\begin{frame}[fragile]
\frametitle{Identify variables involved.}
For Stokes, we have trial variables:
\begin{itemize}
\item fields: scalar $p \in L^{2}(\Omega),$ vector $\vect{u} \in \vect{L}^{2}(\Omega)$, tensor $\NVRtensor{\sigma} \in \vect{L}^{2}(\Omega)$,
\item fluxes: vector $\widehat{\vect{t}}_{n} \in \NVRtensor{H}^{-1/2}(\Gamma_{h})$,
\item traces: vector $\widehat{\vect{u}} \in \vect{H}^{1/2}(\Gamma_{h})$.
\end{itemize}

\begin{lstlisting}
VarFactory varFactory;
VarPtr p = varFactory.fieldVar("p");
VarPtr u1 = varFactory.fieldVar("u_1");
VarPtr u2 = varFactory.fieldVar("u_2");
VarPtr sigma11 = varFactory.fieldVar("\\sigma_11"); 
VarPtr sigma12 = varFactory.fieldVar("\\sigma_12"); 
VarPtr sigma21 = varFactory.fieldVar("\\sigma_21"); 
VarPtr sigma22 = varFactory.fieldVar("\\sigma_22"); 
VarPtr t1n = varFactory.fluxVar("\\hat{t}_1n"); 
VarPtr t2n = varFactory.fluxVar("\\hat{t}_2n"); 
VarPtr u1 = varFactory.traceVar("\\hat{u}_1");
VarPtr u2 = varFactory.traceVar("\\hat{u}_2");
\end{lstlisting}

\end{frame}

\begin{frame}[fragile]
\frametitle{Camellia: Rapid Specification of Inner Products}
Suppose we have a problem whose graph norm is
\[
\NVRnorm{(v,\vect{q})}_{\rm graph}^{2} = \NVRnorm{v}^{2} + \NVRnorm{\vect{q}}^{2} + \NVRnorm{\NVRpd{v}{x} - \NVRpd{v}{y} + \NVRdiv \vect{q}}^{2}.
\]
\vspace{2mm}
To specify this in Camellia, simply do:
\begin{lstlisting}
VarFactory varFactory;
VarPtr v = varFactory.testVar("v", HGRAD);
VarPtr q = varFactory.testVar("q", HDIV);
IPPtr ip = Teuchos::rcp( new IP);
ip->addTerm(v);
ip->addTerm(q);
ip->addTerm(v->dx() - v->dy() + q->div());
\end{lstlisting}

The bilinear form can be specified similarly.

\end{frame}

\begin{frame}
\frametitle{Camellia: Stiffness Matrix Timing Test}
\begin{itemize}
\item local stiffness matrix computation is \pecosbold{embarrassingly parallel}
\item minimize assembly costs: \pecosbold{spatially local} mesh partitioning
\item timing tests on Lonestar: solve convection-dominated diffusion
\end{itemize}
\vspace{-5mm}
\begin{center}
\includegraphics[scale = 0.40]{../figs/scalingFigs/bar_ref3.pdf}\hspace{1cm}
\end{center}
\begin{itemize}
\vspace{-5mm}
\item collaborators working on parallel solvers: Kyungjoo Kim (shared memory architecture),  Maciej Paszynski (distributed memory)
\end{itemize}

\end{frame}

%===============================================================================
% THANK YOU
%===============================================================================
\begin{frame}
\frametitle{}
\begin{block}{}
\center{Thank you!} \\
\center{Questions?}\\
\end{block}
\vspace{1cm}
\center{For more info:\\
nroberts@ices.utexas.edu}\\
\end{frame}

\bibliographystyle{plain}
{\scriptsize
\bibliography{../DPG}
}
 
\end{document}
